\section{Grundlagen}
Dieses Kapitel erläutert die technologischen Grundlagen und Werkzeuge, die für die Durchführung dieser Studienarbeit relevant sind. Zunächst wird die Distributionsplattform Anaconda vorgestellt, gefolgt von einer detaillierten Betrachtung der interaktiven Entwicklungsumgebung Jupyter. Anschließend erfolgt eine Einführung in die Zielsprache TypeScript sowie der Vergleich mit Python. Den Abschluss bildet die Beschreibung der spezifischen Setup- und Toolchain-Konfiguration.


\subsection{Anaconda}
Anaconda ist eine weit verbreitete Open-Source-Distribution für die Programmiersprachen Python und R, die speziell für Aufgaben in den Bereichen Data Science, maschinelles Lernen und wissenschaftliches Rechnen entwickelt wurde. Sie vereinfacht das Paketmanagement und die Bereitstellung von Softwareumgebungen erheblich, was insbesondere in komplexen wissenschaftlichen Projekten von Vorteil ist. Kernkomponente der Distribution ist der Paketmanager \texttt{conda}. Dieser ermöglicht es, Softwarepakete und deren Abhängigkeiten plattformübergreifend zu installieren, zu aktualisieren und zu verwalten. Ein wesentliches Merkmal von Anaconda ist die Fähigkeit, isolierte Umgebungen (Environments) zu erstellen. Diese Isolation verhindert Konflikte zwischen verschiedenen Versionen von Bibliotheken, die für unterschiedliche Projekte benötigt werden könnten.\\ \\
Für die vorliegende Arbeit dient Anaconda als primäre Basisinfrastruktur. Obwohl das Ziel der Portierung auf TypeScript liegt, wird Anaconda ebenfalls verwendet, um die ursprünglichen Python-Notebooks auszuführen und zu validieren. Darüber hinaus fungiert das Anaconda Prompt (bzw. das Terminal) als Schnittstelle für systemweite Befehle. Hierüber erfolgen auch die notwendigen Installationen für das TypeScript-Ökosystem, wie beispielsweise die Ausführung von \texttt{npm install}, um Node.js-Pakete global oder lokal verfügbar zu machen. Die Verwaltung der Node.js-Laufzeitumgebung kann ebenfalls in die bestehende Infrastruktur integriert werden, was Anaconda zu einem vielseitigen Werkzeug macht.


\subsection{Interaktive Programmierung mit Jupyter}
Das Jupyter Notebook ist eine webbasierte, interaktive Entwicklungsumgebung, die es ermöglicht, Dokumente zu erstellen, die sowohl ausführbaren Code als auch formatierten Text, mathematische Formeln und Visualisierungen enthalten. Der Name \enquote{Jupyter} ist ein Akronym, das sich ursprünglich aus den Kernsprachen Julia, Python und R ableitete, mittlerweile aber eine Vielzahl weiterer Programmiersprachen unterstützt.\\ \\
Die Architektur von Jupyter basiert auf einem Client-Server-Modell, das die Trennung zwischen der Benutzeroberfläche und der Code-Ausführung ermöglicht.
\begin{itemize}
	\item Das Frontend (Client): Der Benutzer interagiert über einen Webbrowser mit dem Notebook. Die Oberfläche erlaubt die Eingabe von Code in Zellen sowie die Gestaltung von Textbereichen mittels Markdown.
	\item Der Notebook-Server: Dieser Server verwaltet die Dateien und vermittelt die Kommunikation zwischen dem Frontend und dem ausführenden Kernel.
	\item Das Dateiformat (.ipynb): Jupyter Notebooks werden als \ac{JSON}-Dateien mit der Endung \texttt{.ipynb} gespeichert. Diese strukturierte Textdatei enthält den gesamten Inhalt des Notebooks, Code, Markdown-Text, Metadaten sowie die Ausgabeartefakte (z. B. Konsolenausgaben oder Grafiken), in einem standardisierten Format. Dies erleichtert die Versionierung und den Austausch von wissenschaftlichen Ergebnissen.
\end{itemize}
Ein zentrales Element der Jupyter-Architektur ist der Kernel. Der Kernel ist ein separater Prozess, der für die Ausführung des Codes zuständig ist, den der Benutzer im Notebook eingibt. Wenn eine Code-Zelle ausgeführt wird, sendet das Frontend den Code an den Kernel. Dieser führt die Anweisungen aus und sendet das Ergebnis zurück an das Frontend zur Darstellung.
Für diese Studienarbeit ist das Kernel-Konzept von entscheidender Bedeutung:\\
Während die ursprünglichen Vorlesungsunterlagen auf dem Standard-Python-Kernel (ipykernel) basieren, erfordert die Portierung auf TypeScript die Integration eines dedizierten TypeScript-Kernels (tslab). Diese Architektur ermöglicht den Wechsel der Programmiersprache innerhalb derselben Benutzeroberfläche durch bloßen Austausch des Backends. Die technische Integration des TypeScript-Kernels wird in Abschnitt \hyperref[sec:setup]{Setup und Toolchain} detailliert beschrieben.




\subsection{TypeScript im Data-Science-Kontext}
Während Python als De-facto-Standard im Bereich Data Science gilt, gewinnt das TypeScript-Ökosystem zunehmend an Bedeutung, insbesondere wenn Modelle direkt in Webanwendungen integriert oder visualisiert werden sollen. 

\subsubsection{Herausforderungen und Bibliotheks-Ökosystem}
Eine wesentliche Herausforderung bei der Migration von Python zu TypeScript ist die Diskrepanz in der Standardbibliothek. Während Python mit Modulen wie \texttt{heapq} oder einer nativen Unterstützung für komplexe Mengenoperationen und Tupel ausgestattet ist, fehlen diese Konstrukte in der Standard-Laufzeitumgebung von TypeScript oft oder sind für wissenschaftliche Zwecke nicht performant genug implementiert. Zudem existiert zwar ein reiches Ökosystem (npm), dieses ist jedoch stark fragmentiert, im Gegensatz zu den monolithischen Standard-Paketen in Python wie beispielsweise NumPy. \\
Um die Funktionalität der Python-Vorlesungsunterlagen zu replizieren, wird in dieser Arbeit auf eine Auswahl spezialisierter Bibliotheken zurückgegriffen:

\begin{description}
	\item[Mathematik und Logik:] \hfill \\
	Für numerische Stabilität und symbolische Berechnungen kommen \texttt{mathjs} (umfassende Mathematik-Bibliothek) und \texttt{fraction.js} (für rationale Zahlen zur Vermeidung von Fließkommafehlern) zum Einsatz. Für Constraint-Solving-Probleme und logische Beweise werden der \texttt{logic-solver} sowie der \texttt{z3-solver} (ein Theorem-Prover von Microsoft Research) eingebunden.
	
	\item[Datenstrukturen:] \hfill \\
	Da TypeScript keine native Priority Queue besitzt (äquivalent zu Pythons \texttt{heapq}), wird \texttt{heap-js} verwendet. Eine Besonderheit stellt die Bibliothek \texttt{recursive-set} dar. Native JavaScript \texttt{Sets} prüfen Objekte auf Referenzgleichheit, nicht auf Wertegleichheit. Das bedeutet, zwei inhaltlich identische Objekte werden als unterschiedlich behandelt. \texttt{recursive-set} ermöglicht echte Mengenoperationen auf Basis der tiefen Gleichheit (Deep Equality).
	
	\item[Visualisierung und Medien:] \hfill \\
	Zur Generierung von Graphen wird \texttt{@hpcc-js/wasm} genutzt. Dies ist eine WebAssembly-Portierung von Graphviz, die das Rendern von DOT-Sprache ermöglicht, ohne dass eine lokale \texttt{graphviz.exe}-Installation auf dem Host-System notwendig ist – ein entscheidender Vorteil für die Portabilität des Notebooks. Für die Bildmanipulation wird \texttt{pngjs} verwendet.
	
	\item[Type Definitions:] \hfill \\
	Ein Spezifikum der TypeScript-Entwicklung ist die Notwendigkeit separater Typ-Definitionen für Bibliotheken, die ursprünglich in reinem JavaScript geschrieben wurden. So muss beispielsweise für die Bildverarbeitung zusätzlich das Paket \texttt{@types/pngjs} als \textit{Dev-Dependency} installiert werden, um die statische Typprüfung zu ermöglichen.
\end{description}

\textbf{@TobiKnight}
\subsubsection{Syntaxvergleich und Typisierungskonzepte}
Der fundamentale Unterschied zwischen Python und TypeScript liegt im Zeitpunkt und der Strenge der Typprüfung. \\


\noindent \textbf{Statisch vs.\ Dynamisch / MyPy vs.\ TypeScript Compiler}\\
Um diesen Unterschied zu verstehen, ist es hilfreich, zunächst zu klären, was ein \textit{Datentyp} ist.
Jeder Wert, den ein Programm verarbeitet, hat eine bestimmte Art: Eine ganze Zahl wie \texttt{42}
verhält sich anders als ein Text wie \texttt{"Hallo"} oder ein Wahrheitswert wie \texttt{True}.
Diese verschiedenen Arten von Werten nennt man \textit{Datentypen}.
Der grundlegende Unterschied zwischen Python und TypeScript liegt darin, \textit{wann} und
\textit{ob überhaupt} geprüft wird, ob diese Typen im Programm korrekt verwendet werden.\\


\noindent \textbf{Dynamische Typisierung in Python}
\\
Python verwendet standardmäßig \textit{dynamische Typisierung}.
Das bedeutet, dass der Typ einer Variable nicht im Voraus festgelegt werden muss.
Python ermittelt den Typ erst dann, wenn das Programm tatsächlich ausgeführt wird, also
zur sogenannten \textit{Laufzeit}.
Als Analogie lässt sich ein Rucksack vorstellen, in den man beliebige Dinge (Datentypen) packen kann:
Mal liegt ein Buch darin, mal eine Flasche Wasser.
Erst wenn man etwas herausnehmen will und feststellt, dass es nicht passt, gibt es ein Problem.\\
In Python ist es daher ohne Weiteres möglich, einer Variable zunächst eine Zahl
und später einen Text zuzuweisen:

\Needspace{6\baselineskip}
\begin{minted}{python}
 x = 42        # x hat den Typ int (ganze Zahl)
 x = "Hallo"   # x hat jetzt den Typ str (Text)  Python erlaubt dies
 # Der Fehler tritt erst beim Ausführen der nächsten Zeile auf:
 x + 1
\end{minted}
\captionof{listing}{Dynamische Typisierung in Python}

Den Fehler in der letzten Zeile würde Python erst bemerken, wenn diese Zeile tatsächlich ausgeführt wird.
In einem großen Programm mit vielen Zeilen Code kann das dazu führen, dass Fehler lange
unentdeckt bleiben, zum Beispiel dann, wenn eine bestimmte Stelle im Code nur selten erreicht wird.

\subsubsection{Syntaxvergleich und Typisierungskonzepte}
Der fundamentale Unterschied zwischen Python und TypeScript liegt im Zeitpunkt und der Strenge der Typprüfung.\\
Um diesen Unterschied zu verstehen, ist es hilfreich, zunächst zu klären, was eine \textit{Variable} und was ein \textit{Datentyp} ist.
In der Programmierung speichert man Daten in Variablen. Man kann sich eine Variable wie eine Kiste vorstellen: Der \textit{Name} der Variable (z.\,B. \texttt{x}) steht als Etikett außen auf der Kiste, und der eigentliche \textit{Wert} (z.\,B. \texttt{42}) liegt darin. Jeder dieser Werte hat zudem eine bestimmte Art: Eine ganze Zahl wie \texttt{42} verhält sich anders als ein Text wie \texttt{"Hallo"} oder ein Wahrheitswert wie \texttt{True}. Diese verschiedenen Arten von Werten nennt man \textit{Datentypen}.

Der grundlegende Unterschied zwischen Python und TypeScript liegt darin, \textit{wann} und \textit{ob überhaupt} geprüft wird, ob diese Typen im Programm korrekt verwendet werden.\\


\textbf{Dynamische Typisierung in Python}\\
Python verwendet standardmäßig \textit{dynamische Typisierung}.
Das bedeutet, dass der Typ einer Variable nicht im Voraus festgelegt werden muss.
Python ermittelt den Typ erst dann, wenn das Programm tatsächlich ausgeführt wird, also zur sogenannten \textit{Laufzeit}.
Als Analogie lässt sich eine unbeschriftete Schublade vorstellen. Zuerst legt man einen Hammer (eine Zahl) hinein. Später räumt man die Schublade aus und legt stattdessen einen Pinsel (einen Text) hinein, die Schublade akzeptiert beides problemlos. Der Fehler fällt erst in dem Moment auf, in dem man blind in die Schublade greift, um einen Nagel in die Wand zu schlagen: Man merkt erst beim Zugriff auf die Schublade (Variable), dass man mit einem Pinsel nicht hämmern kann.


In Python ist es daher ohne Weiteres möglich, einer Variable zunächst eine Zahl und später einen Text zuzuweisen:

\Needspace{6\baselineskip}
\begin{minted}{python}
 x = 42        # x hat den Typ int (ganze Zahl)
 x = "Hallo"   # x hat jetzt den Typ str (Text) -- Python erlaubt dies
 # Der Fehler tritt erst beim Ausfuehren der naechsten Zeile auf:
 x + 1
\end{minted}
\captionof{listing}{Dynamische Typisierung in Python}

Den Fehler in der letzten Zeile würde Python erst bemerken, wenn diese Zeile tatsächlich ausgeführt wird.
In einem großen Programm mit vielen Zeilen Code kann das dazu führen, dass Fehler lange unentdeckt bleiben, zum Beispiel dann, wenn eine bestimmte Stelle im Code nur selten erreicht wird.\\

\textbf{Statische Typisierung in TypeScript}

TypeScript hingegen verwendet \textit{statische Typisierung}.
Der Typ einer Variable muss entweder explizit angegeben werden oder wird vom Compiler automatisch erkannt. Diese automatische Erkennung nennt man \textit{Typinferenz} (abgeleitet von \textit{inferieren}, also schlussfolgern).
Entscheidend ist, dass TypeScript alle Typen bereits \textit{vor} der Ausführung des Programms überprüft.
Dieser Vorgang heißt \textit{Kompilierung}, und das Programm, das diese Prüfung durchführt, nennt man \textit{TypeScript-Compiler}.

Der TypeScript-Compiler liest den gesamten Code durch und prüft, ob alle Typangaben widerspruchsfrei sind.
Findet er einen Fehler, bricht er ab und gibt eine Fehlermeldung aus. Das Programm wird also gar nicht erst gestartet.

Das folgende Beispiel zeigt, wie TypeScript denselben Fehler verhindert. Eine Methode um in TypeScript eine neue Variable anzulegen, ist die Verwendung des Schlüsselworts \texttt{let}. Es signalisiert dem Compiler: Hier wird eine neue Variable erstellt.  Dahinter folgt der Name der Variable, ein Doppelpunkt und der gewünschte Datentyp gefolgt von einem Gleichzeichen gefolgt von dem zuzuweisenden Wert. \\

\Needspace{6\baselineskip}
\begin{minted}{typescript}
 let x: number = 42;
 // Die folgende Zeile wuerde vom Compiler sofort rot markiert werden:
 x = "Hallo";
 // Fehlermeldung: Type 'string' is not assignable to type 'number'.
\end{minted}
\captionof{listing}{Statische Typisierung in TypeScript}

Der Compiler würde bei der Zuweisung des Textes sofort einen Fehler melden, da in eine \texttt{number}-Variable kein \texttt{string} (Text) gelegt werden darf.
Der Vorteil ist, dass solche Fehler bereits beim Schreiben des Codes und nicht erst beim Ausführen gefunden werden.\\

\textbf{MyPy als optionale Typprüfung für Python}

Da die dynamische Typisierung in größeren Projekten fehleranfällig sein kann, wurde für Python das Werkzeug \textit{MyPy} entwickelt.
MyPy ist ein \textit{statischer Typprüfer} für Python: Es analysiert den Code auf Typfehler, indem es ihn lediglich durchliest, ohne das Programm tatsächlich zu starten. Anders als der TypeScript-Compiler, der den Code nach der Prüfung auch übersetzt, ist MyPy ein reines Kontrollwerkzeug und belässt den Python-Code unverändert.


Ein wesentlicher Unterschied zum TypeScript-Compiler ist, dass MyPy vollständig \textit{optional} ist.
Um MyPy nutzen zu können, müssen im Code freiwillige Typangaben vorhanden sein, sogenannte \textit{Type Hints} (Typhinweise).
Diese Typhinweise haben keinen Einfluss auf das Verhalten des Programms, Python ignoriert sie bei der Ausführung vollständig.

Das folgende Beispiel verdeutlicht dies mit einer einfachen Variable:

\Needspace{6\baselineskip}
\begin{minted}{python}
 # Die Variable wird mit dem Typhinweis ': int' (ganze Zahl) versehen
 alter: int = 20
 # MyPy meldet hier VOR dem Start einen Fehler, da "Zwanzig" ein Text ist.
 alter = "Zwanzig"
\end{minted}
\captionof{listing}{Typhinweise in Python mit MyPy-Prüfung}


\textbf{Bedeutung für die Übersetzung}

Für die Übersetzung von den Python-Notebooks in TypeScript-Notebooks ergibt sich aus dem Unterschied der Typprüfung eine wichtige Herausforderung.
Die Python-Notebooks können in drei Zuständen vorliegen:

\begin{enumerate}
	\item \textbf{Vollständig ungetypter Code:} Es sind keine Typangaben vorhanden.
	Die Typen müssen also beim Übersetzen ermittelt und ergänzt werden.
	\item \textbf{Teilweise getypter Code:} Nur manche Stellen enthalten Typhinweise.
	Hier müssen einerseits die fehlenden Typen für TypeScript ermittelt und ergänzt werden. Andererseits können die bereits vorhandenen Typangaben oft nicht einfach eins zu eins übernommen werden, sondern erfordern ebenfalls Anpassungen.
	\item \textbf{Vollständig mit MyPy annotierter Code:} Alle relevanten Stellen enthalten Typhinweise.
	Diese können zwar oft direkt in TypeScript-Typen übertragen werden, jedoch müssen selbst hier häufig noch Präzisierungen vorgenommen werden. Der Grund dafür ist, dass das Typsystem von TypeScript in manchen Bereichen anders strukturiert oder strenger ist als das von MyPy.
	
\end{enumerate}

\subsection{Setup und Toolchain}
\label{sec:setup}
Die Reproduzierbarkeit der Ergebnisse und die Stabilität der Entwicklungsumgebung sind essenziell für die Validierung der portierten Artefakte. Dieses Unterkapitel beschreibt die Architektur der verwendeten Toolchain und dokumentiert die notwendigen Konfigurationsschritte, wie sie in der \texttt{\href{https://github.com/zLeoll/TS-Translation-for-TheoInf/blob/main/README.md}{README.md}} der Arbeit definiert sind.

\subsubsection{Architektur der Laufzeitumgebung}
Die technische Infrastruktur basiert auf einer hybriden Architektur, die das Python-zentrierte Ökosystem von Anaconda mit der JavaScript-basierten Node.js-Laufzeitumgebung verbindet. Diese Trennung ist notwendig, da Jupyter ursprünglich für Python konzipiert wurde, TypeScript jedoch eine Node.js-Laufzeit für die Transpilierung und Ausführung benötigt.\\
Die Toolchain gliedert sich in drei Schichten:
\begin{enumerate}
	\item \textbf{Basisinfrastruktur (Anaconda):} Bereitstellung von Python, Jupyter (\texttt{nbclassic}) und Verwaltung der Systempfade.
	\item \textbf{Laufzeitumgebung (Node.js):} Ausführungsumgebung für den TypeScript-Compiler und die genutzten Bibliotheken.
	\item \textbf{Kernel-Layer (tslab):} Eine Middleware, die das Jupyter-Messaging-Protokoll implementiert und TypeScript-Codezellen an die Node.js-Laufzeit weiterleitet.
\end{enumerate}

\subsubsection{Implementierung der Umgebung}
Die Einrichtung der Umgebung erfolgt deklarativ über die Kommandozeile der Anaconda-Distribution. Zunächst wird die Laufzeitumgebung durch die Installation von Node.js und dem klassischen Jupyter-Frontend vorbereitet. Dies schafft die Voraussetzungen für die Integration nicht-nativer Kernel.
\Needspace{6\baselineskip}
\begin{minted}{typescript}
	conda activate base
	conda install nodejs
	conda install -c conda-forge nbclassic
\end{minted}
\captionof{listing}{Bereitstellung der Laufzeitumgebung}

Für die Ausführung von TypeScript innerhalb der Notebooks wird der Kernel \texttt{tslab} verwendet. Im Gegensatz zu reinen JavaScript-Kerneln (wie IJavascript) ermöglicht \texttt{tslab} eine direkte Typ-Prüfung innerhalb der Zellen, was für die didaktische Zielsetzung der Arbeit, die Vermittlung von Typsicherheit, unabdingbar ist.

\Needspace{6\baselineskip}
\begin{minted}{typescript}
	tslab install
	npm install -g tslab
	npm install typescript
\end{minted}
\captionof{listing}{Integration des TypeScript-Kernels}

Die fachlichen Anforderungen der Vorlesungen erfordern spezifische Bibliotheken, die über den Paketmanager \texttt{npm} bezogen werden. Hierbei wird strikt zwischen Laufzeit-Abhängigkeiten (z.\,B. \texttt{mathjs} für Numerik, \texttt{recursive-set} für Mengenlehre) und Entwicklungs-Abhängigkeiten (z.\,B. \texttt{@types/pngjs}) unterschieden. Letztere sind notwendig, um für ursprünglich in JavaScript geschriebene Bibliotheken eine statische Typprüfung zu ermöglichen.

\Needspace{6\baselineskip}
\begin{minted}{typescript}
	// Mathematische und logische Kernkomponenten
	npm install fraction.js mathjs logic-solver z3-solver
	
	// Datenstrukturen und Mengenoperationen
	npm install heap-js recursive-set@8.0.0
	
	// Visualisierung und Medienverarbeitung
	npm install @hpcc-js/wasm pngjs
	npm i --save-dev @types/pngjs
\end{minted}
\captionof{listing}{Installation der Projekt-Abhängigkeiten}

\subsubsection{Validierung der Toolchain}
Um sicherzustellen, dass die Interaktion zwischen dem Jupyter-Frontend, dem tslab-Kernel und den installierten Node-Modulen korrekt funktioniert, wurde ein dediziertes Testverfahren etabliert. Das Referenz-Notebook \href{https://github.com/zLeoll/TS-Translation-for-TheoInf/blob/main/Algorithms/TypeScript/Chapter-02/Test-Libraries.ipynb}{Test-Libraries.ipynb} führt einen systematischen Import-Test aller Komponenten durch. Dieses Notebook verifiziert, dass:
\begin{itemize}
	\item der Kernel korrekt registriert ist und TypeScript-Syntax verarbeitet,
	\item alle externen Module im Pfad der Laufzeitumgebung gefunden werden,
	\item die Typdefinitionen korrekt aufgelöst werden.
\end{itemize}
Erst nach erfolgreicher Ausführung dieses Notebooks gilt die Umgebung als valide für die Bearbeitung der Vorlesungsunterlagen.